\documentclass[11pt]{article}
\usepackage[T1]{fontenc}
\usepackage[utf8]{inputenc}
\usepackage{amsmath,amssymb,amsthm,geometry,hyperref}
\geometry{margin=1in}
\setlength{\emergencystretch}{3em}
\newtheorem{proposition}{Proposition}
\newtheorem{theorem}{Theorem}
\newcommand{\ddc}{dd^c}
\newcommand{\PSH}{\operatorname{PSH}}
\newcommand{\NP}{\mathrm{NP}}
\newcommand{\Ca}{\mathrm{Cao}}
\newcommand{\vol}{\operatorname{vol}}
\newcommand{\Area}{\operatorname{Area}}
\newcommand{\B}{\mathcal B}
\title{A divisorial partial construction for zero-volume transcendental Okounkov bodies}
\author{Unverified draft for original P15.3}
\date{4 October 2026}
\begin{document}
\maketitle
\begin{abstract}
We prove a restricted current-valuation construction for rational divisorial singularities with smooth semipositive remainder. It retains the ordinary envelope type in this domain and is independent of the stated positive perturbations. A nonprojective transcendental product example has class volume zero, intrinsic body dimension two, and a computed unnormalized two-dimensional mass depending on the prescribed logarithmic coefficient. We do not construct a theory for arbitrary nonanalytic P-types or prove a general numerical-dimension theorem. Thus unchanged original P15.3 remains open in this draft. Execution, compilation, and hashes are not mathematical verification.
\end{abstract}

\section{Exact source scope}
The only mathematical sources read are the four supplied primary caches: Darvas--Reboulet--Witt Nystr\"om--Xia--Zhang (DRWNXZ), arXiv:2309.07584v3, 5 January 2026; Xia, arXiv:2112.04290v2, 13 April 2024; Cao, arXiv:1210.5692v1, 21 October 2012; and Choi--Hyun--Park--Won (CHPW), arXiv:1508.03922v5, whose arXiv header says 5 February 2018 and internal manuscript date says 24 September 2018. Source locations below are one-based indices of the supplied UTF-8 text split on literal LF, retaining form feeds and the trailing empty line. Full literal passages, relevant proof bodies, versions, locations, hashes, and application checks are preserved in \path{source-extracts/source-provenance.json}.

DRWNXZ Definition 3.2 and the pseudoeffective definition (lines 561--691) use successive generic Lelong numbers of analytic currents along a globally smooth connected complete flag. Their Theorem 7.6 and Corollary 7.7 (1389--1415) prove ambient volume equals class volume divided by $n!$. The nonbig proof compares with positive perturbations; it does not identify lower-dimensional numerical data or incorporate a prescribed potential. Definition 5.1 and Example 5.3 (852--894), however, impose bounded singularity order on analytic K\"ahler currents. We use this genuinely partial construction at positive perturbations.

Xia's section spaces $H^0(X,L^k\otimes\mathcal I(k\phi))$ (73--99), Proposition 5.8 (1181--1185), and section 5.6 (1623--1637) make his body I-invariant. His definitions distinguish bounded singularity type, the ordinary $P$-envelope, the mass-sensitive $C$-envelope, and the I-envelope (371--429). Proposition 2.10 identifies $C$ with $P$ only when the actual top nonpluripolar mass is positive. Theorem 2.19 (490--494) computes the section volume as the actual mass of $P[\phi]_{\mathcal I}$ divided by $n!$, not unconditionally as the actual mass of $\phi$. Positive class volume, positive ideal-model volume, and positive actual mass are different assertions.

Section 5.6 defines a limit using rational positive ample perturbations. Conjecture 5.29 explicitly states $\dim\Delta(\theta,\phi)=\operatorname{nd}(\theta,\phi)$ and refers to Cao Definition 4. No proof is supplied there. Xia's construction is projective and line-bundle based; we do not apply it to the transcendental example.

Cao Definitions 2 and 3, with Proposition 2.2 and its proof (162--279), use the absolutely continuous parts of quasi-equisingular analytic approximants to define a cohomological product of general positive currents. Definition 4 (284--289) defines the dimension of a line bundle with metric from that product. This is not the definition of the actual nonpluripolar product of the original current. We write $d_{\Ca}(T)$ for the largest nonzero degree of the current product in Definition 3; we do not treat a transcendental class as a line bundle.

CHPW Theorem B/Corollary 3.22 (173--185, 1071--1085) requires a smooth projective variety, a pseudoeffective real divisor, and a flag containing a positive volume subvariety of the numerical Iitaka dimension. Its lower-dimensional mass is augmented restricted volume along that subvariety. The relevant proofs are Theorem 3.20 and Proposition 3.21 (919--1070). Remark 3.24 (1090--1096) allows strict dimension inequality for other flags. These algebraic theorems are not applied below. Our flag calculation is direct. These source limitations are not an impossibility proof for P15.3.

\section{Divisorial construction and all claimed independence}
We use $\ddc=i\partial\bar\partial/(2\pi)$, so $\ddc\log|z|^2$ is the unit divisor current. Fix a compact K\"ahler $n$-fold $M$ with a globally smooth connected complete flag $Y_\bullet$, a smooth prime divisor $F$, a rational $a>0$, and $D=aF$. A smooth Hermitian metric on $\mathcal O(F)$ has curvature $e$. Suppose $\beta$ is a smooth closed semipositive form and set
\[
 \phi=a\log|s_F|_h^2,\qquad \theta=\beta+ae,\qquad
 T=\theta+\ddc\phi=\beta+[D],\qquad \alpha=[\theta].
\]
Assume $\vol(\alpha)=0$. A specified flag, including its specified modification if needed, is part of the reference data. For a K\"ahler form $\Omega$ define
\[
 \B_Y(\alpha,\phi)=\bigcap_{\varepsilon>0}
 \Delta_Y\bigl(\alpha+\varepsilon[\Omega];T+\varepsilon\Omega\bigr).
 \tag{2.1}
\]
At positive $\varepsilon$ this is DRWNXZ Definition 5.1: valuations are those of K\"ahler currents with gentle analytic singularities whose potentials are at most $\phi$ plus a bounded constant. The reference current $T+\varepsilon\Omega$ is K\"ahler and gently analytic because $D$ is a rational smooth divisor. No artificial convex shape or appended potential label is involved.

\begin{proposition}
The set (2.1) is nonempty, compact, convex, and
\[
 \B_Y(\alpha,\phi)=\nu_Y([D])+\Delta_Y([\beta]). \tag{2.2}
\]
It is independent of the K\"ahler perturbation form. It retains the ordinary $P$-envelope type of $\phi$ in this rational divisorial bounded-type domain. Its dimension is $\dim\Delta_Y([\beta])$; no general numerical-dimension equality is asserted.
\end{proposition}
\begin{proof}
At positive $\varepsilon$, Example 5.3 of DRWNXZ applies with smooth residual $\beta+\varepsilon\Omega$. We spell out its two inclusions. A residual gently analytic K\"ahler current $S$ gives the admissible current $S+[D]$, with valuation $\nu_Y(S)+\nu_Y([D])$. Conversely an admissible analytic K\"ahler current $G$ has generic multiplicity at least $a$ on $F$. DRWNXZ Lemma 2.5 subtracts $[D]$ and retains positive analytic singularities. In fact the residual is K\"ahler: subtract from $G-\delta\Omega$ for a sufficiently small $\delta>0$. Its valuation belongs to the residual class body by Definition 3.2 and Lemma 3.1. Consequently
\[
 \Delta_Y(\alpha+\varepsilon[\Omega];T+\varepsilon\Omega)
 =\nu_Y([D])+\Delta_Y([\beta]+\varepsilon[\Omega]). \tag{2.3}
\]
Taking intersections gives (2.2) by the actual pseudoeffective class-body definition. Compactness and convexity follow from that construction. The smooth residual at every $\varepsilon$ shows that $\nu_Y([D])$ belongs to every member, so the intersection is nonempty. This proof makes no nonanalytic P-preserving approximation claim.

Adding a smooth positive form preserves all valuation coordinates. Thus positive smooth differences of big classes give inclusions of their bodies. For K\"ahler forms $\Omega,\Omega'$ choose $0<c<C$ with $c\Omega<\Omega'<C\Omega$ pointwise. The corresponding residual bodies sandwich each other after rescaling $\varepsilon$. Their intersections therefore agree; no zero-volume comparison is used.

More generally, suppose $\rho_j$ are smooth K\"ahler forms with norm tending to zero relative to a fixed $\Omega$. For suitable $\delta_j\to0$,
\[
 \Delta_Y([\beta])\subseteq\Delta_Y([\beta]+[\rho_j])
 \subseteq\Delta_Y([\beta]+\delta_j[\Omega]).
\]
Nested nonempty compact sets converge in Hausdorff distance to their intersection: otherwise points remaining a fixed distance from the intersection admit a convergent subsequence in a fixed member; its limit lies in every member, a contradiction. The sandwich proves convergence along these positive paths to (2.2). Rational cofinal $\varepsilon$ give the same intersection. We do not claim convergence for arbitrary nonpositive perturbations or for positive forms that fail to tend to zero.
\end{proof}

We next prove the retained P assertion, independently of any I-envelope identification. With fixed $\theta$ write
\[
 P_\theta[\phi]=\left(\sup\{v\in\PSH(M,\theta):v\leq0,
                  \ v\leq\phi+C_v\}\right)^*.
\]
For each competitor its generic divisorial Lelong number is at least $a$. Subtracting $[D]$ leaves a positive residual with a $\beta$-psh potential $w=v-\phi$, extended across $F$ as its upper semicontinuous potential. Off $F$, $v\leq0$ gives the bound $w\leq-\phi$, independent of $C_v$. Locally use $F=\{z_1=0\}$, $\phi=a\log|z_1|^2+b$ with $b$ smooth, and a local smooth potential $r$ of $\beta$. On transverse circles $|z_1|=r_0$, with other coordinates in a smaller polydisc, $w+r$ has a uniform upper bound by $-\phi+r$. On each transverse disc it is subharmonic or identically $-\infty$. The maximum principle extends this bound to the disc interior. A finite cover of $F$ and the compact complement give $w\leq K$ uniformly for all competitors. Hence $v\leq\phi+K$, including after upper semicontinuous regularization, by the local logarithmic bound. The competitor $\phi-\sup_M\phi$ supplies a lower bound. We have proved
\[
                 P_\theta[\phi]=\phi+O(1). \tag{2.4}
\]
Boundedly equivalent inputs have identical competitor families, and (2.4) gives $P_\theta[P_\theta[\phi]]=P_\theta[\phi]$ in this domain. Changing the metric on $\mathcal O(F)$ changes $\phi$ by a bounded smooth function and changes neither the constraint nor the envelope data nor the body. This proves retention of the ordinary analytic P-type. It makes no reconstruction claim for arbitrary types from one flag.

\subsection{Actual versus Cao products in the stated domain}
Off $F$ the current $T$ equals the smooth form $\beta$. The actual nonpluripolar product is local there and charges no pluripolar set, so
\[
                    \langle T^k\rangle_{\NP}=\beta^k. \tag{2.5}
\]
The constant sequence $\phi_j=\phi$ is a quasi-equisingular analytic approximation in Cao Definition 2: convergence and positivity are immediate, the singularity order is unchanged, and the required ideal inclusion follows by normalizing $\phi\leq0$, since integrability against the larger exponent implies integrability against the smaller one. The absolutely continuous part of $T$ is $\beta$. Thus Cao Definition 3 gives $[\beta^k]$. Equality with the cohomology of (2.5) has been proved in this particular domain, not assumed generally.

This cohomological data persists for the bounded type, including the P-envelope in (2.4). Indeed after subtracting $[D]$ a current of that bounded type has residual $\beta+\ddc b$, where $b$ is globally bounded. Bounded Bedford--Taylor products satisfy
\[
 (\beta+\ddc b)^k-\beta^k
 =\ddc\left(b\sum_{j=0}^{k-1}(\beta+\ddc b)^j\wedge\beta^{k-1-j}\right).
 \tag{2.6}
\]
This is the telescoping identity with the defining local bounded-potential operation and multiplication by closed products. These products charge no pluripolar sets, agree off $F$ with the actual nonpluripolar products, and give their extension across $F$. Equation (2.6) proves equality of cohomology classes, not equality of the measures of different representatives.

\section{A transcendental example and a complete restricted numerical proof}
Let $S$ be a compact K\"ahler surface and $\kappa$ a K\"ahler class outside the real span of divisor classes, with fixed K\"ahler representative $\kappa_0$. For existence, one may take a complex two-torus with no nonzero integral $(1,1)$-classes: choose a linear complex structure on $\mathbb R^4$ outside the countably many proper real-analytic loci where a fixed nonzero integral alternating two-form is invariant. Each is proper by varying the complex structure to fail invariance. Baire's theorem supplies such a structure, and a constant Hermitian form makes $\mathbb R^4/\mathbb Z^4$ K\"ahler. This is an existence argument, not a computed sample.

Choose $p\in S$ and $\pi:\widehat S=\operatorname{Bl}_pS\to S$, with exceptional divisor $E$. In a coordinate ball take a smooth cutoff $\chi$ supported there and equal to one near $p$, and put
\[
 f=\chi\log(|z_1|^2+|z_2|^2),\qquad u_a=af.
\]
For sufficiently small $a>0$, $\kappa_0+\ddc u_a\geq0$. Near $p$ the logarithmic Hessian is semipositive, while on the compact cutoff annulus the negative Hessian is bounded by a fixed multiple of $\kappa_0$. Choose rational $a>0$ below that bound. Then
\[
 \pi^*(\kappa_0+\ddc u_a)=\beta_a+a[E],\qquad
 [\beta_a]=\kappa_a:=\pi^*\kappa-a[E], \tag{3.1}
\]
where $\beta_a$ is smooth K\"ahler. To prove this, in the chart $z_1=w,z_2=wt$, the potential is $a\log|w|^2+a\log(1+|t|^2)$. Subtracting $a[E]$ leaves
\[
                  \pi^*\kappa_0+a\ddc\log(1+|t|^2).
\]
On $E$ the first form is positive normally and the second tangentially. In a neighborhood both are semipositive with complementary kernels; the other chart is identical. Small $a$ gives strict positivity on the cutoff annulus, and outside it $\pi$ is an isomorphism with the original positive form. This proves smoothness and strict positivity, rather than assuming an unproved numerical bound.

Let $C$ be any connected compact Riemann surface, $c_0\in C$, and $q\in E$. Start on $X_0=S\times C$, with $\alpha_0=p_S^*\kappa$ and $\phi_0=p_S^*u_a$. Use the fixed modification
\[
 \Pi=\pi\times\operatorname{id}_C:X=\widehat S\times C\longrightarrow X_0.
\]
Its center $\{p\}\times C$ is smooth. Put $\alpha=\Pi^*\alpha_0$, $\theta=p_{\widehat S}^*\pi^*\kappa_0$, $\phi_a=\Pi^*\phi_0$, $\beta=p_{\widehat S}^*\beta_a$, and $D=a(E\times C)$. Then $\theta+\ddc\phi_a=\beta+[D]$. Use the globally smooth connected flag
\[
 X\supset Y_1=\widehat S\times\{c_0\}
  \supset Y_2=E\times\{c_0\}\supset Y_3=\{(q,c_0)\}. \tag{3.2}
\]
This is the geometric current-valuation flag of the transcendental theorem, not an assertion about the transcendence degree of a nonprojective meromorphic function field.

The example is genuinely transcendental. If $\alpha_0$ were in the real span of divisor classes, restriction to $S\times\{c\}$ would put $\kappa$ in that span. For $\alpha$ the same argument uses restriction to $\widehat S\times\{c\}$ followed by $\pi_*$; divisor classes push to divisor classes or zero. The indicated torus example is nonprojective. The specified modification supplies the complete flag.

The setup of DRWNXZ section 2 initially fixes a big class, so we verify its psh pullback argument in the present nonbig case. Given $v\in\PSH(X,\Pi^*\theta_0)$, off the exceptional image the function $w=v\circ\Pi^{-1}$ is $\theta_0$-psh and bounded above because $v$ is quasi-psh on a compact space. Locally add a smooth potential of $\theta_0$; the bounded-above psh function extends across the analytic exceptional image by its upper semicontinuous extension, exactly the removable-singularity argument in the proof of Lemma 2.1. Its pullback and $v$ agree off an analytic set and hence almost everywhere, so their canonical psh representatives agree everywhere by submean regularization. This proves surjectivity without a bigness hypothesis. Injectivity and pointwise order preservation follow from surjectivity of $\Pi$.

The proved order isomorphism preserves $v\leq0$ and both bounded-order inequalities. Envelopes are least psh upper bounds of their competitor families, so it commutes with the ordinary P-envelope for the pulled-back form and zero obstacle. Thus the retained type is the pullback of the prescribed original point-logarithmic P-type. No independence of unrelated modifications or flags is asserted.

\begin{theorem}[Restricted product theorem]
For these data,
\[
 \B_Y(\alpha,\phi_a)=
 \{(0,a+x,y):(x,y)\in\Delta_{(E,q)}(\kappa_a)\}. \tag{3.3}
\]
Its affine carrier is $\{x_1=0\}$ with standard $(x_2,x_3)$ coordinate area, and its intrinsic dimension is exactly two. Its unnormalized lower-dimensional mass is
\[
 2!\Area(\B_Y)=\int_{\widehat S}\beta_a^2
              =\int_S\kappa_0^2-a^2>0. \tag{3.4}
\]
Its ambient volume, class volume, and actual top nonpluripolar mass are all zero. The actual degree-two nonpluripolar cohomological product and Cao's current product are $[\beta^2]\ne0$, and their largest nonzero degree is two. The nef class $\alpha$ also has numerical dimension two. For fixed $p$, the body recovers the prescribed rational logarithmic coefficient via $\min x_2=a$.
\end{theorem}
\begin{proof}
First, $\theta$ is smooth semipositive. There is no K\"ahler current $K$ in $\alpha$: if $K\geq\delta\Omega$ for a K\"ahler $\Omega$, pairing with $\theta^2$ would give
\[
 0=\int_X K\wedge\theta^2\geq\delta\int_X\Omega\wedge\theta^2>0.
\]
Here $\theta^3=0$ because it is pulled back from a surface, while $\theta^2$ is nonzero and positive on an open set. Thus $\alpha$ is nonbig and its class volume is zero by the published definition. The same argument works on $X_0$. This is not an inference from the mass of one potential.

Choose a K\"ahler form $\eta$ on $\widehat S$ and a positive form $\sigma$ on $C$ with $\int_C\sigma=1$. Put $\Omega_0=p_{\widehat S}^*\eta+p_C^*\sigma$. For every $\varepsilon>0$ we claim
\[
 \Delta_Y([\beta]+\varepsilon[\Omega_0])
 =[0,\varepsilon]\times\Delta_{(E,q)}(\kappa_a+\varepsilon[\eta]). \tag{3.5}
\]
For the upper inclusion, take a positive globally analytic current $G$ in the indicated class and let $t=\nu(G,Y_1)$. Subtract $t[Y_1]$, leaving a positive analytic current. Its restriction to a generic curve $\{z\}\times C$ is defined: choose $z$ so this curve is not contained in the analytic singularity locus. The degree of that restriction is $\varepsilon-t$, since the product fiber divisor meets that curve once. Hence $0\leq t\leq\varepsilon$. After the first valuation subtraction, restriction to $Y_1$ is positive and analytic near the flag point; its class is $\kappa_a+\varepsilon[\eta]$, because the normal bundle of the product fiber is trivial and the curve pullback restricts trivially. Its two remaining coordinates lie in the surface body by DRWNXZ Definition 3.2/Lemma 3.1. Passing to closure gives the upper inclusion.

For the converse, take a surface gently analytic K\"ahler current $R$ in $\kappa_a+\varepsilon[\eta]$. For rational $0\leq t<\varepsilon$ the current
\[
 p_{\widehat S}^*R+p_C^*((\varepsilon-t)\sigma+t[c_0])
\]
is K\"ahler, has rational gentle analytic singularities, and has valuation $(t,\nu_{(E,q)}(R))$. Strict positivity holds in both product directions. Rational $t$ are dense in $[0,\varepsilon]$ with the last endpoint obtained in closure; such surface valuations are dense in the surface body. This proves the lower inclusion and (3.5).

The surface class $\kappa_a$ is K\"ahler, hence big. DRWNXZ Proposition 3.3(3) gives
\[
 \bigcap_{\varepsilon>0}\Delta_{(E,q)}(\kappa_a+\varepsilon[\eta])
 =\Delta_{(E,q)}(\kappa_a).
\]
Intersecting (3.5) yields $\Delta_Y([\beta])=\{0\}\times\Delta_{(E,q)}(\kappa_a)$. The fixed divisor has valuation $(0,a,0)$: it does not contain $Y_1$, restricts to $aE$, and leaves zero residual after the second subtraction. Proposition 1 proves (3.3). The surface body has nonempty interior because its class is K\"ahler, so the dimension is exactly two. Neither zero ambient volume nor an unproved dimension conjecture has been used to reach this conclusion.

DRWNXZ Theorem 7.6 applied on the surface gives its area as $\vol(\kappa_a)/2!$. A K\"ahler class has volume its self-intersection: bounded truncations have that total top mass by Stokes, all nonpluripolar masses are bounded above by it, and a smooth representative attains it. Since $\pi^*\kappa\cdot E=0$ by contraction and $E^2=-1$ from the normal bundle $\mathcal O_E(-1)$, expansion gives (3.4).

Equation (2.5) gives the actual products $\beta^k$. Degree three vanishes by the product pullback, while degree two is nonzero by pairing with $\Omega_0$. The constant analytic approximation gives the same Cao cohomological products, so $d_{\Ca}(T)=2$. The smooth nef representative $\theta$ likewise has largest nonzero intersection degree two. Bounded-type and P-envelope cohomological invariance are proved in (2.6).

The geometric dimensional carrier is the surface $Y_1$ and its residual mass is $\int_{\widehat S}\beta_a^2$. Equivalently, for the explicitly named degree-one transverse scale,
\[
 \int_X\beta^2\wedge\Omega_0=\int_{\widehat S}\beta_a^2.
\]
Changing $\eta$ or the degree-one $\sigma$ does not change this number. A different reference form of fiber degree $m$ would give $m$ times the number; we make no assertion that changing this scale is irrelevant. No zero mass is normalized to a probability.

The surface body contains $(0,0)$ from a smooth representative and all its coordinates are nonnegative. Consequently (3.3) contains $(0,a,0)$ and has minimum second coordinate $a$. Distinct small rational $a$ therefore give distinct bodies in the same original class $\alpha_0$, with lower mass $\kappa^2-a^2$. This recovers the coefficient of the prescribed P-type for fixed center $p$ through an actual valuation constraint. It does not recover arbitrary supports or nonanalytic types from one flag. All perturbation independence is exactly the positive-form and positive-path independence proved in Proposition 1 on the specified model.
\end{proof}

\section{An exact obstruction to analytic domination of arbitrary P-data}
The analytic restriction above is substantive. Take $M=\mathbb P^1\times C$, for a compact connected curve $C$, and $\theta=p_1^*\omega_{\mathrm{FS}}$, where $\int_{\mathbb P^1}\omega_{\mathrm{FS}}=1$. Its class is pseudoeffective and nonbig: a K\"ahler current in $[\theta]$, paired with $\theta$, contradicts $\theta^2=0$ and $\int_M\Omega\wedge\theta>0$. Thus its class volume is zero.

Choose distinct dense points $p_j\in\mathbb P^1$ and $a_j=2^{-j-2}$ for $j\geq1$, so $A=\sum_j a_j=1/4$. Let $g_j\leq0$ be the logarithm of the squared chordal distance to $p_j$. Then $\ddc g_j=[p_j]-\omega_{\mathrm{FS}}$, its pole has coefficient one, and its $L^1(\omega_{\mathrm{FS}})$ norm is independent of $p_j$ by unitary invariance. Define
\[
 u=\sum_{j\geq1}a_jp_1^*g_j,\qquad F_j=\{p_j\}\times C.
\]
The partial sums decrease, are $\theta$-psh, and have summable $L^1$ norms. Their limit is therefore $\theta$-psh and not identically $-\infty$. Distributional convergence gives the explicit positive current
\[
 \theta+\ddc u=(1-A)\theta+\sum_{j\geq1}a_j[F_j].
\]
The actual top nonpluripolar mass also vanishes: every bounded truncation $\max(u,-k)$ is pulled back from the curve, so its local top wedge on the surface is zero. This is a separate computation from class volume zero.

Put $\psi=P_\theta[u]$ with zero obstacle. The competitor $u\leq0$ ensures a nontrivial envelope. Since $u\leq a_jp_1^*g_j$ for every $j$, its competitor family is contained in that for $P_\theta[a_jp_1^*g_j]$. The single smooth-divisor envelope proof in Section 2 applies with residual $(1-a_j)\theta$, giving
\[
 \psi\leq P_\theta[a_jp_1^*g_j]\leq a_jp_1^*g_j+K_j.
\]
Its uniform transverse logarithmic bound survives upper semicontinuous regularization. Thus $\psi=-\infty$ on each $F_j$, whose union is dense.

\begin{proposition}[Empty analytic-domination competitors]
For this prescribed envelope $\psi$, no nontrivial quasi-psh $v\leq\psi+C_v$ has analytic singularities on any nonempty open neighborhood. In particular, for every K\"ahler $\Omega$ and every $\varepsilon>0$, no such competitor in $\PSH(M,\theta+\varepsilon\Omega)$ is analytic even near a chosen flag center.
\end{proposition}
\begin{proof}
Such a $v$ is $-\infty$ along every $F_j$. If it were analytic on an open neighborhood, in a smaller coordinate neighborhood its $-\infty$ locus would be a proper closed zero set of finitely many holomorphic functions, or empty. It contains the dense union of the fibers intersecting that neighborhood, hence is the entire neighborhood. Every generator would vanish identically and $v$ would be identically $-\infty$ there, contradicting local integrability of a nontrivial quasi-psh potential. The argument is independent of the background class of $v$.
\end{proof}
This defeats the specified recipe of valuing analytic currents more singular than a fixed P-envelope at each positive perturbation and intersecting. It does not exclude another genuine geometric or valuation realization for the original problem, and is not an impossibility inference from an unsuccessful search.

\section{Remaining original obligations and precise status}
The domain of the proposition is rational divisorial analytic data with smooth semipositive residual. The dimension theorem is restricted further to the constructed product. The unchanged original problem includes arbitrary compact K\"ahler manifolds, arbitrary zero-volume pseudoeffective classes, and prescribed P-singularity data. The following are not proved here:
\begin{enumerate}
\item A genuine geometric or current-valuation body for arbitrary nonanalytic P-envelope types without replacing them by I-equivalence. Proposition 2 shows that simply retaining the analytic-domination rule can leave every positive-perturbation competitor set empty.
\item P-preserving analytic approximation, nonemptiness, compactness, or limiting behavior for that general construction. I-equisingularity does not establish these P assertions.
\item A general lower-dimensional numerical theorem for the retained data, with the actual product or other invariant precisely identified. Neither Cao's product nor Xia's conjectured equality is made into such a theorem here.
\item The exact general flag hypotheses realizing numerical dimension. The product calculation proves only the specified flag (3.2).
\item All general modification relations or independence claims when a complete smooth flag is absent. One specified modification and flag are retained here.
\item Any additional reference-scale, carrier, and perturbation-independence assertions needed in that generality. The present sandwiches use a fixed divisor and smooth residual.
\end{enumerate}
There is no impossibility proof for the original demand. A source conjecture, a restricted theorem, or the obstruction to the specified recipe does not give such a proof. P15.1 and P15.2 are outside the assignment; no energy or Monge--Amp\`ere operator is constructed. The required status is partial and the candidate flag is false.

\section{Documentary evidence and rerun instructions}
The frozen input files are \path{input.json}, its byte-identical copy \path{input-original.json}, the exact decoded assignment in \path{assignment-original.txt}, and the frozen supplied task object in \path{task-original.json}. The last two standalone files did not initially exist; their derivation and initial hashes are in \path{input-preservation.json}. Subsequent scripts do not rewrite them.

Run \texttt{python3 prepare\_sources.py} to regenerate literal source records from the four allowed caches; \texttt{python3 build\_paper.py} for two actual builds, each with preserved stdout, stderr, and TeX log; and \texttt{python3 verify\_artifacts.py} for documentary consistency checks. Tool versions and OS/kernel are in \path{environment.json}. Commands, return codes, and diagnostics are preserved. The parallel prose argument is \path{P15.3-argument.md}. The required five-field \path{status.json} records the exact partial result and concrete gaps. No exhaustive numerical experiment or independent mathematical review is claimed.

\begin{thebibliography}{4}
\bibitem{D} T. Darvas, R. Reboulet, D. Witt Nystr\"om, M. Xia, K. Zhang, \emph{Transcendental Okounkov bodies}, arXiv:2309.07584v3, 2026.
\bibitem{X} M. Xia, \emph{Partial Okounkov bodies and Duistermaat--Heckman measures of non-Archimedean metrics}, arXiv:2112.04290v2, 2024.
\bibitem{C} J. Cao, \emph{Numerical dimension and a Kawamata--Viehweg--Nadel-type vanishing theorem on compact K\"ahler manifolds}, arXiv:1210.5692v1, 2012.
\bibitem{H} S. R. Choi, Y. Hyun, J. Park, J. Won, \emph{Okounkov bodies associated to pseudoeffective divisors}, arXiv:1508.03922v5, 2018.
\end{thebibliography}
\end{document}
